\documentclass{ProblemSetCUNY}
\usepackage{bussproofs} \EnableBpAbbreviations
\usepackage{mathrsfs}
\usepackage{stix}
\usepackage{pgffor}
\usepackage{listings}

\usepackage{tabstackengine}
\TABstackMath
\TABbinary

\AssignmentNumber{2}%
\CourseName{Automation Engineering}
\CourseNumber{39531}
\InstructorName{Alex Washburn}
\DueDateYear{2026}
\DueDateMonth{02}
\DueDateDay{11}


%\newenvironment{\ContactListing}[args]{begdef}{enddef}

\newcommand{\Bool}{\ensuremath{\mathtt{Bool}}\xspace}
%\newcommand{\Nats}{\ensuremath{\mathbb{N}}\xspace}
%\newcommand{\Ints}{\ensuremath{\mathbb{Z}}\xspace}
%\newcommand{\Rats}{\ensuremath{\mathbb{Q}}\xspace}
%\newcommand{\Reals}{\ensuremath{\mathbb{R}}\xspace}
%\newcommand{\ProofLabelOne}[1]{\RightLabel{\ensuremath{\Parens{\mathrm{\textsc{#1}}}}\xspace}}
\newcommand{\ProofLabelTwo}[2]{\RightLabel{\ensuremath{\Parens{\mathrm{\textsc{#1}}}\Parens{\mathrm{#2}}}\xspace}}

\newcommand{\LogicCase}[1]{\item \textbf{Case}~#1:\newline}

\newcommand{\DerivationFormat}[2]{\hfill\qquad#2\textsc{#1}}
\newcommand{\Derivation}[2]{\DerivationFormat{#1}{#2~\ensuremath{\ruledelayed}~}}

\newcommand{\AxiomCL}[1]{\DerivationFormat{Axiom~#1}{}}
\newcommand{\Factivity}{\DerivationFormat{Factivity}{}}
\newcommand{\Given}{\DerivationFormat{Hypothesis}{}}
\newcommand{\Contrapositive}{\DerivationFormat{Contrapositive}{}}
\newcommand{\DoubleNegation}{\DerivationFormat{Double Negation}{}}
\newcommand{\DeMorgans}{\DerivationFormat{DeMorgans}{}}
\newcommand{\Necitation}[1]{\Derivation{Necitation}{#1}}
\newcommand{\Distribution}[1]{\Derivation{Distribution}{#1}}
\newcommand{\Reflexivity}[1]{\Derivation{Reflexivity}{#1}}
\newcommand{\ConDef}[1]{\Derivation{Connective Definition}{#1}}
\newcommand{\Contradiction}[1]{\ensuremath{\neg \phi \land \phi}\Derivation{Contradiction}{\ensuremath{\phi = #1}}}
\newcommand{\DedThm}[1]{\Derivation{Deduction Theorem}{#1}}
\newcommand{\ModPon}[2]{\Derivation{Modus Ponens}{#1, #2}}
\newcommand{\Syllogism}[2]{\Derivation{Syllogism}{#1, #2}}
\newcommand{\Theorem}[2]{\DerivationFormat{Theorem~--~Lecture~#1,~Page~#2}{}}

\newcommand{\R}[2]{\ensuremath{\World{#1}\mathrel{R}\World{#2}}\xspace}
\newcommand{\Tr}[2]{\ensuremath{\tau\Parens{#1,\,#2}\xspace}}
\newcommand{\T}[1]{\ensuremath{\tau\Parens{#1}\xspace}}
\newcommand{\Prob}[1]{\ensuremath{P\Parens{#1}}\xspace}
\newcommand{\ProbGiven}[2]{\ensuremath{P\Parens{#1~\,|\,~#2}}\xspace}
\newcommand{\RandVar}[1]{\ensuremath{\mathbf{#1}}\xspace}
\newcommand{\DistMap}[2]{~#1&\mapsto~~#2\\}
\newcommand{\DistCondition}[2]{~#1&\text{\normalfont iff}~~#2\\}


\newcommand{\One}{\ensuremath{\textbf{1}}\xspace}

\newcommand{\Cons}{%
\mathrel{:\!:}%
}

\newcommand{\ProblemItem}[1]{%
\hspace{\parindent}\textbullet\;\;#1%
}

\begin{document}
\CoverPage%


%
% Mapping
% Configuration updates
% Folding
%


\Problem{%
Kara Vann has signal processing function that they need to write.
They tried writing one big loop, but it got too complicated and Kara Vann got stuck with defective code.
Kara Vann describes the signal processing function to you and asks for your help.
They need to eventually produce the following function:
$$\mathtt{makeSignal} \colon \mathtt{Stream}\;\,\Ints \to \mathtt{Stream}\;\,\Bool$$
}%
%
\SubProblem{Write a function which ensures that the value is in the range. $\NumericRange{-10}{10}$.\newline%
The function should have the type signature: $\mathtt{clamp} \colon \Ints \to \Ints$}
\vfill
\SubProblem{Write a function which divides an even number by 2, subtracts 1 from a negative odd number, and adds 1 to a positive odd number.\newline%
The function should have the type signature: $\mathtt{perturb} \colon \Ints \to \Ints$}
\vfill
\SubProblem{Write a function which outputs the minimum distance from the input to either -5 or 5.\newline%
The function should have the type signature: $\mathtt{measure} \colon \Ints \to \Nats$}
\vfill
\SubProblem{Kara Vann says that $\mathtt{makeSignal}$ should do the following, in order:\newline%
\ProblemItem{$\mathtt{clamp}$ the stream values}\newline%
\ProblemItem{then $\mathtt{perturb}$ the values}\newline%
\ProblemItem{then \emph{add 1} to the values}\newline%
\ProblemItem{then $\mathtt{perturb}$ the values}\newline%
\ProblemItem{then $\mathtt{measure}$ the values}\newline%
\ProblemItem{then $\mathtt{perturb}$ the values}\newline%
\ProblemItem{then \emph{halve} the values}\newline%
\ProblemItem{finally, output if the value is even}\newline%
Write the function $\mathtt{makeSignal}$ by composing together the functions above and anonymous functions.}
\vfill

\Problem{%
Nora Mally has defined a user configuration data-type:\\%
\begin{equation*}
\mathtt{Config}\; \Coloneq \begin{cases}
\mathtt{DefaultInfo}\\
\mathtt{UserDefined}\\
\quad\;\Tuple{\mathtt{darkmode} \colon \Bool,\, \mathtt{columnWidth} \colon \Nats,\, \mathtt{fontScaling} \colon \Rats }
\end{cases}
\end{equation*}\\[5mm]
Nora Mally wants to always initialize their application with the $\mathtt{DefaultInfo}$ for all users.
Then whenever a user changes one of the labeled fields of their configuration, the $\mathtt{Config}$ is updated.
Instead of writing a boring ``setter'' method for each field of $\mathtt{Config}$, Nora Mally wants to use abstractions to simplify their code.
They figure if they add or remove fields later, they will have a simpler time maintaining the code.\\[5mm]%
\textit{Help Nora Mally write methods so that $\mathtt{Config}$ satisfies the constraint classes below.}
}

\SubProblem{Define a $\mathtt{Semigroup}$ instance for $\mathtt{Config}$.\newline%
\ProblemItem{Write a function $\oplus \colon \mathtt{Config} \to \mathtt{Config} \to \mathtt{Config}$.}\newline%
\ProblemItem{Show that your $\oplus$ function satisfies associativity.}%
}
\vfill

\SubProblem{Define a $\mathtt{Monoid}$ instance for $\mathtt{Config}$.\newline%
\ProblemItem{Define a value $\mathtt{zed} \colon \mathtt{Config}$.}\newline%
\ProblemItem{Show that your $\mathtt{zed}$ value is an identity element for your $\oplus$ function.%
}}
\vfill


\Problem{%
Ed Vance is planning to transform some data.
They want to write some \emph{generic} analysis functions over $\mathtt{Folable}$ structures.
Sometimes an empty structure cannot return a result that Ed Vance wants from the generic function, so they defined the following data-type to represent ``missingness'' or ``existence'':\\%
\begin{equation*}
\mathtt{Optional}\;\alpha \Coloneq \begin{cases}
\mathtt{None}\\
\mathtt{Some}\;\alpha
\end{cases}
\end{equation*}\\[5mm]
\textit{Help Ed Vance write their desired functions below by using $$\Parens{\mathtt{Foldable}\; t } \Rightarrow \mathtt{fold} \colon \Parens{\alpha \to \beta \to \beta} \to \beta \to t\;\alpha \to \beta$$}
}
\SubProblem{$\Parens{\mathtt{Foldable}\; t } \Rightarrow \mathtt{length} \colon t\;\alpha \to \Nats$}
\vfill
\SubProblem{$\Parens{\mathtt{Foldable}\; t,\, \mathtt{Equatable}\; \alpha } \Rightarrow \mathtt{contains} \colon t\;\alpha \to \alpha \to \Bool$}
\vfill
\SubProblem{$\Parens{\mathtt{Foldable}\; t,\, \mathtt{Orderable}\; \alpha } \Rightarrow \mathtt{maximum} \colon t\;\alpha \to \mathtt{Optional}\;\alpha$}
\vfill
\SubProblem{$\Parens{\mathtt{Foldable}\; t } \Rightarrow \mathtt{first} \colon t\;\alpha \to \mathtt{Optional}\;\alpha$}
\vfill
\SubProblem{$\Parens{\mathtt{Foldable}\; t } \Rightarrow \mathtt{last} \colon t\;\alpha \to \mathtt{Optional}\;\alpha$}
\vfill


\end{document}
